\begin{remarks}[5.7.3]
\label{IV.2.s5.7.3-fr}
\begin{enumerate}
\item[(i)] Recordemos que siempre se tiene
$\operatorname{prof}(\mathscr F_x)\leq\dim(\mathscr F_x)$ (\textbf 0, 16.4.5.1)
si $\mathscr F_x\neq0$. Decir que $\mathscr F$ posee la propiedad $(S_k)$
significa, pues, que se tiene $\operatorname{prof}(\mathscr F_x)\geq k$ \emph{salvo}
en los puntos $x\in X$ tales que $\dim(\mathscr F_x)<k$, y que en estos
últimos puntos se tiene
$\dim(\mathscr F_x)=\operatorname{prof}(\mathscr F_x)$, es decir
(\textbf 0, 16.5.1), que $\mathscr F$ es un $\mathcal O_X$-módulo \emph{de
Cohen-Macaulay} en estos puntos; se observará que, en los puntos en que
$\dim(\mathscr F_x)=k$, se tiene $\operatorname{prof}(\mathscr F_x)=k$,

\oldpage[IV-2]{104}

por lo que $\mathscr F$ sigue siendo un $\mathcal O_X$-módulo de
Cohen-Macaulay en estos puntos. Decir que $\mathscr F$ posee la propiedad
$(S_k)$ \emph{para todo} $k$ significa, pues, que $\mathscr F$ es un
$\mathcal O_X$-\emph{módulo de Cohen-Macaulay}. Es claro que, para $k'\geq k$, la
propiedad $(S_{k'})$ implica $(S_k)$; para $k\leq0$, \emph{todo}
$\mathcal O_X$-módulo coherente posee la propiedad $(S_k)$.

\item[(ii)] Para verificar la condición (5.7.2.1), es posible limitarse al caso
en que $x\in\operatorname{Supp}(\mathscr F)$; en caso contrario se tiene, en
efecto, $\dim(\mathscr F_x)=-\infty$ (\textbf 0, 14.1.2).

\item[(iii)] Si $X=\operatorname{Spec}(A)$, donde $A$ es un anillo
noetheriano, y $\mathscr F=\widetilde M$, donde $M$ es un $A$-módulo de tipo
finito, se dice que $M$ posee la propiedad $(S_k)$ si $\mathscr F$ posee esta
propiedad. Para un preesquema localmente noetheriano cualquiera $X$, decir que
$\mathscr F$ posee la propiedad $(S_k)$ \emph{en un punto} $x\in X$ significa, pues,
que, si se pone $Y=\operatorname{Spec}(\mathcal O_x)$, el
$\mathcal O_Y$-módulo $\widetilde{\mathscr F}_x$ posee la propiedad $(S_k)$;
se observará que la condición (5.7.2.1), en general, \emph{no implica} (5.7.2.2) para
toda generización $x'$ de $x$, puesto que no se tiene ninguna relación de
desigualdad entre $\operatorname{prof}(\mathscr F_x)$ y
$\operatorname{prof}(\mathscr F_{x'})$ (\textbf 0, 16.4.6).

\item[(iv)] Se sigue inmediatamente de la definición que, si $\mathscr F$
verifica $(S_k)$ en un punto $x$, verifica también $(S_k)$ en todo punto $x'$
que sea una \emph{generización} de $x$.

\item[(v)] La propiedad $(S_k)$ es especialmente importante para $k=1$ y
$k=2$; fue introducida para $k=2$ por Serre, para expresar su criterio de
normalidad (cf. (5.8.5)).

\item[(vi)] Sean $X$ un preesquema localmente noetheriano, $Y$ un
subpreesquema cerrado de $X$, $j:Y\to X$ la inyección canónica, y
$\mathscr G$ un $\mathcal O_Y$-módulo coherente. Es claro que para todo
$x\in Y$ se tiene
$\dim(\mathscr G_x)=\dim((j_*(\mathscr G))_x)$ y
$\operatorname{prof}(\mathscr G_x)=\operatorname{prof}((j_*(\mathscr G))_x)$,
de donde
$\operatorname{coprof}(\mathscr G_x)=\operatorname{coprof}((j_*(\mathscr G))_x)$.
Para que $\mathscr G$ verifique $(S_k)$, es necesario y suficiente que
$j_*(\mathscr G)$ verifique $(S_k)$.
\end{enumerate}
\end{remarks}

\begin{proposition}[5.7.4]
\label{IV.2.s5.7.4-fr}
Sean $X$ un preesquema localmente noetheriano y $\mathscr F$ un
$\mathcal O_X$-módulo coherente. Pongamos
$S=\operatorname{Supp}(\mathscr F)$ y, para todo $n\geq0$, designemos por
$Z_n$ el conjunto de los $z\in X$ tales que
$\operatorname{coprof}(\mathscr F_x)>n$, de modo que $Z_n\subset S$.
\begin{enumerate}
\item[\textup{(i)}] Para que $\mathscr F$ verifique la propiedad $(S_k)$, es necesario
y suficiente que, para todo $n\geq0$, se tenga
\begin{equation}
\operatorname{codim}(Z_n,S)>n+k.
\tag{5.7.4.1}\label{IV.2.s5.7.4.1-fr}
\end{equation}
\item[\textup{(ii)}] Supongamos además que los $Z_n$ son cerrados en $X$. Entonces,
para que $\mathscr F$ verifique la propiedad $(S_k)$ en un punto $x\in S$, es
necesario y suficiente que, para todo $n\geq0$, se tenga
\begin{equation}
\operatorname{codim}_x(Z_n,S)>n+k.
\tag{5.7.4.2}\label{IV.2.s5.7.4.2-fr}
\end{equation}
\end{enumerate}
\end{proposition}

(i) En efecto, por definición (5.1.3), se tiene
$\operatorname{codim}(Z_n,S)=\inf_{z\in Z_n}(\dim(\mathcal O_{S,z}))$, y la
desigualdad (5.7.4.1) significa, pues, (5.1.12.2), que, para todo $z\in X$ y
todo $n\geq0$, la relación
\[
\dim(\mathscr F_z)-\operatorname{prof}(\mathscr F_z)\geq n+1
\]
implica la relación
\[
\dim(\mathscr F_z)\geq n+k+1
\]
Pero, si se pone $a=\dim(\mathscr F_z)$,
$b=\operatorname{prof}(\mathscr F_z)$, se tiene $b\leq a$, y decir que, para
todo $n\geq0$, la relación $b\leq a-n-1$ implica $k\leq a-n-1$ equivale a
decir que $b\geq\inf(k,a)$, de donde resulta la proposición.

\oldpage[IV-2]{105}

(ii) El razonamiento es el mismo que en (i), salvo que hay que limitarse a
los $z\in X$ que son \emph{generizaciones} de $x$, y tener en cuenta (5.1.3.2).
